\section{项目需求}

\subsection{项目背景}
大禾众邦（厦门）智能科技股份有限公司是国内冷弯薄壁型钢建筑领域具有软硬件独立研发能力的领军企业，长期专注于轻钢龙骨设备、C/Z型钢成型设备、智能生产线及配套设计软件的研发与推广。公司产品与解决方案覆盖建筑工业化、装配式建筑、旅居用房、宿居用房、楼宇及库房等多种应用场景，并已参与《装配式冷弯薄壁型钢结构施工及验收规程》等全国性标准的制定工作，积累了丰富的行业经验与工程数据。

随着装配式建筑市场对结构安全性、设计效率与成本控制的要求不断提高，传统冷弯薄壁型钢设计流程中“设计—建模—结构验算—加工输出”环节相互割裂、数据无法贯通的问题日益凸显。为充分发挥企业在设备、软件与工程实践方面的综合优势，亟需建设一套面向冷弯薄壁型钢结构的“一体化设计平台”，实现从建筑设计、有限元分析到规范校核的数据闭环与流程贯通。

\subsection{平台建设目标}
本项目拟为大禾众邦开发面向冷弯薄壁型钢结构的一体化设计平台中的核心模块——有限元分析模块。平台整体流程为：

\begin{center}
\textbf{建筑设计} \quad$\longrightarrow$\quad \textbf{有限元分析} \quad$\longrightarrow$\quad \textbf{规范校核}
\end{center}

其中，有限元分析模块作为连接几何模型与规范验算的关键环节，需要实现以下能力：

\begin{enumerate}[label=\arabic*., leftmargin=2em]
    \item 接收并解析建筑三维模型与构件参数，自动生成可用于计算的几何与材料数据；
    \item 对冷弯薄壁型钢构件及整体结构进行静力、稳定与模态分析；
    \item 输出应力、变形、屈曲模态等关键结果，为后续规范校核提供可靠依据；
    \item 与现有 BuildingMaster（BM）设计软件及生产控制系统进行数据接口对接，形成“设计—分析—生产”一体化闭环。
\end{enumerate}

\subsection{有限元分析的必要性、难点与痛点}

\subsubsection{必要性}
冷弯薄壁型钢构件具有壁薄、截面形式复杂、局部屈曲与畸变屈曲敏感等显著特点，传统基于简化公式的手算或经验设计方法难以准确反映其真实力学行为。有限元分析能够：

\begin{itemize}[leftmargin=2em]
    \item 精确模拟开口薄壁截面在轴压、弯曲、剪切及组合作用下的应力分布与变形特征；
    \item 考虑材料非线性、几何非线性及初始缺陷对构件承载力的影响；
    \item 对整体结构的稳定性、连接节点性能及抗震性能进行定量评估；
    \item 为优化截面形式、减少用钢量、提升结构安全冗余度提供科学依据。
\end{itemize}

\subsubsection{难点与痛点}
尽管有限元分析在学术界已较为成熟，但在冷弯薄壁型钢建筑企业的工程应用中仍面临诸多挑战：

\begin{itemize}[leftmargin=2em]
    \item \textbf{建模效率低：}现有通用有限元软件（如ANSYS、ABAQUS）操作复杂，需要专业 CAE 工程师手动完成几何清理、网格划分、边界条件施加等步骤，难以适应快速设计需求。
    \item \textbf{截面与规范脱节：}通用软件通常不内置中国冷弯薄壁型钢相关规范（如《冷弯薄壁型钢结构技术规范》GB 50018 等），设计人员需在后处理阶段人工提取结果并自行校核，容易出错且效率低下。
    \item \textbf{初始缺陷与残余应力处理困难：}薄壁构件对初始几何缺陷、残余应力及边界条件极为敏感，通用软件中相关参数的设置缺乏针对性和工程经验数据。
    \item \textbf{数据孤岛：}设计软件、分析软件与生产控制系统之间缺乏统一数据接口，导致设计变更后需反复导出、转换、校验，严重影响设计—制造协同效率。
\end{itemize}

\subsection{项目需求概述}
基于上述背景，本项目拟开发一套面向冷弯薄壁型钢结构、贴合大禾众邦业务需求的有限元分析模块，核心需求如下：

\begin{enumerate}[label=\arabic*., leftmargin=2em]
    \item \textbf{自动化建模：}支持由建筑信息模型（BIM）或设计软件直接导入构件几何、材料与荷载信息，实现有限元模型的快速自动生成。
    \item \textbf{多类型分析：}覆盖静力分析、稳定性分析、模态分析等常用工况，满足轻钢建筑结构验算的基本需求。
    \item \textbf{规范对接：}分析结果能够直接支撑后续规范校核环节，为截面强度、构件稳定及整体承载力验算提供数据接口。
    \item \textbf{轻量化与可集成：}模块应具备良好的可扩展性，能够嵌入大禾众邦现有设计软件体系，支持本地化部署与云端协同。
\end{enumerate}

通过本项目的实施，有望显著提升大禾众邦在冷弯薄壁型钢建筑领域的设计效率、结构安全评估能力与数字化水平，进一步增强其“设计—分析—生产”一体化解决方案的市场竞争力。

